\documentclass[article, 11pt]{article}

%%%%%%%%%%%%%%%%%%%%%%%%%%%%%%%%%%%%%%%%%%%%%%%%%%%
% Preamble:
%%%%%%%%%%%%%%%%%%%%%%%%%%%%%%%%%%%%%%%%%%%%%%%%%%%
\usepackage{fullpage}
\usepackage{amsfonts}
\usepackage{amsmath}
\usepackage{amsthm}
\usepackage{graphicx}
\usepackage{color}
\usepackage{palatino, url, multicol}
\usepackage{enumerate}
\usepackage{ulem}
\usepackage{hyperref}
\usepackage{verbatim}

\thispagestyle{empty}

%%%%%%%%%%%%%%%%%%%%%%%%%%%%%%%%%%%%%%%%%%%%%%%%%%%%
% Start Document:
%%%%%%%%%%%%%%%%%%%%%%%%%%%%%%%%%%%%%%%%%%%%%%%%%%%%
\begin{document}

% Project Title:
\begin{center}
\boxed{\LARGE{\textbf{\textsc{MINI-PROJECT 4}}}}\\
\Huge{\textbf{\textsc{Differential Equations \& Modeling: The Two-Body Problem}}}\\
\end{center}

\vspace{.6cm}

% Brief Summary:
\noindent \boxed{\textbf{\textsc{Brief Summary}}}

\vspace{-.35cm}
\noindent \hrulefill

This mini-project uses differential equations to derive and simulate the classical two-body problem under Newtonian gravity. You will (1) derive the system of ODEs governing the motion of two masses interacting via an inverse-square gravitational force, (2) derive and interpret the motion of the center of mass, and (3) numerically solve and visualize trajectories for chosen initial conditions. A key goal is connecting the mathematical derivation to a system of nonlinear differential equations for which we can computationally/numerically solve in Python to get an idea of the trajectory of the bodies.  

\vspace{.35cm}

% Goals:
\noindent \boxed{\textbf{\textsc{Goals}}}

\vspace{-.35cm}
\noindent \hrulefill

\begin{itemize}
\item Derive the two-body problem equations of motion from Newton's laws and Newton's law of gravitation.
\item Define the center of mass and formulate how the position vector of the center of mass changes depending on each individual body.
\item Reduce and/or rewrite the model in a convenient first-order ODE system for numerical simulation.
\item Use Python to simulate trajectories and communicate results through clear figures and a short animation/movie.
\end{itemize}

\vspace{.35cm}

% LaTeX Deliverables:
\noindent \boxed{\textbf{\textsc{\LaTeX\, Deliverables}}}

\vspace{-.26cm}
\noindent \hrulefill

\begin{itemize}
\item A short \LaTeX\, report that is clearly written and well-organized.
\item Derive the two-body problem system of equations starting from Newton's laws and the gravitational force.  Include a derivation for the center of mass.  Also, if you want to show that the center of mass has constant velocity, include that as well (not necessary for full credit). 
\item Clearly define all variables and parameters.
\item Include at least one figure with a caption that shows a trajectory of the two-body problem for a set of initial conditions.   The figure must also show the center of massInclude axis labels, a legend, and a title.
\item Provide an interpretation or description of the motion, in particular noting how the center of mass moves. 
\end{itemize}

\vspace{.35cm}

% Python Deliverables:
\noindent \boxed{\textbf{\textsc{Python Deliverables}}}

\vspace{-.35cm}
\noindent \hrulefill

\begin{itemize}
\item A Python script that sets parameters $G$, $m_1$, and $m_2$ and initial conditions (positions and velocities), writes the model as a first-order ODE system (introduce velocities as additional variables), numerically solves the system over a chosen time interval, and produces and saves the required trajectory figure used in the \LaTeX\ write-up.
  
\item The Python script also creates and saves a movie/animation of the motion that shows both bodies and the center of mass. 
\item All plots must have titles, axis labels, legends, and readable formatting.

\end{itemize}

\vspace{.35cm}

% Suggested Analysis Questions (optional):
\noindent \boxed{\textbf{\textsc{Suggested Analysis Questions}}}

\vspace{-.28cm}
\noindent \hrulefill


\begin{itemize}
\item How do different initial velocities change the qualitative behavior?
\item How does changing the mass ratio $m_1/m_2$ affect each body's path relative to the center of mass?
\end{itemize}

\vspace{.35cm}

% AI Disclosure:
\noindent \boxed{\textbf{\textsc{AI Disclosure:}}}

\vspace{-.35cm}
\noindent \hrulefill

If you use AI in any way, you must include an Appendix in the write-up that notes the LLM used and provides a list of the prompts you used. You may use AI for suggestions related to Python code, \LaTeX\, syntax, or other general ideas, but the report should not be copy-and-pasted from AI.

\pagebreak

\end{document}